\subsection{The Pipeline as One Algorithm}
\label{sec:algorithm}

The pipeline is one algorithm that builds a surrogate from a dataset of solution curves. Its inputs are the dataset, its data ceiling per target and its physical domain. The ceiling is a property of the incoming dataset, measured from it in the audit (Sections~\ref{sec:eda-ns} and~\ref{sec:eda-bh}), and not a step of the automated process; the domain is the physical edge of Section~\ref{sec:domain-boundaries}. The same algorithm runs on the two-parameter black-hole map and on the three-parameter neutron-star map.

\begin{algorithm}[!htb]
\caption{Building a surrogate of a compact-object map.}
\label{alg:surrogate}
\begin{algorithmic}[1]
\Require a dataset of solution curves, its data ceiling per target, its physical domain
\Ensure per target, a surrogate with its relative error beside the ceiling, which refuses an input outside the domain
\State prepare the dataset and cut each curve at its stable edge (Section~\ref{sec:preprocessing})
\State split by whole curves, so no curve is in training and test at once
\State choose the representation: input scale and target transform (Section~\ref{sec:representation})
\State screen the model families on the same folds (Section~\ref{sec:families})
\State tune the survivors under one search budget (Section~\ref{sec:optimization})
\State refit the winner with its numerical precision (Section~\ref{sec:precision})
\State guard it with the physical domain (Section~\ref{sec:validity})
\end{algorithmic}
\end{algorithm}
